% Distilled blueprint for Evans, Chapter 9 (Nonvariational Techniques).
% Formal declarations, metadata, and citation identifiers are preserved; exposition and exercises are omitted.

\section{Monotonicity Methods}
\label{sec:monotonicity-methods}

Let $U\subset\R^n$ be bounded with smooth boundary, let $f\in L^2(U)$, and consider
\[
\begin{cases}
-\Div\,\mathbf a(Du)=f & \text{in }U,\\
u=0 & \text{on }\partial U,
\end{cases}
\tag{1}
\]
where $\mathbf a:\R^n\to\R^n$ is smooth.

\begin{definition}[Monotone vector field]
\label{def:monotone-vector-field}
\dcref{ch9:9.1-def-1}
\group{Nonlinear Elliptic PDE}
A \textit{vector field} $\mathbf{a} : \R^n \to \R^n$ is called \textit{monotone} provided
\[
(\mathbf{a}(p) - \mathbf{a}(q)) \cdot (p - q) \geq 0
\tag{3}
\]
for all $p, q \in \R^n$.
\end{definition}

Assume $\mathbf a$ is monotone and, for constants $C,\alpha>0$ and $\beta\geq0$,
\[
|\mathbf a(p)|\leq C(1+|p|),
\tag{4}
\]
\[
\mathbf a(p)\cdot p\geq\alpha|p|^2-\beta
\tag{5}
\]
for all $p\in\R^n$. Let $\{w_k\}_{k=1}^\infty$ be a smooth orthonormal basis of $H_0^1(U)$ for the inner product
$\int_U Du\cdot Dv\,dx$. The Galerkin approximations have the form
\[
u_m=\sum_{k=1}^m d_m^k w_k
\tag{6}
\]
and satisfy
\[
\int_U\mathbf a(Du_m)\cdot Dw_k\,dx=\int_Ufw_k\,dx
\quad(k=1,\ldots,m).
\tag{7}
\]
A weak solution of (1) is a function $u\in H_0^1(U)$ such that
\[
\int_U\mathbf a(Du)\cdot Dv\,dx=\int_Ufv\,dx
\quad(v\in H_0^1(U)).
\tag{12}
\]

\begin{lemma}[Zeros of a vector field]
\label{lem:zeros-of-a-vector-field}
\uses{thm:brouwer-fixed-point-theorem}
\dcref{ch9:9.1-lem-1}
\group{Nonlinear Elliptic PDE}

Assume the continuous function $\mathbf{v} : \R^n \to \R^n$ satisfies
\[
\mathbf{v}(x) \cdot x \geq 0 \quad \text{if } |x| = r,
\tag{8}
\]
for some $r > 0$. Then there exists a point $x \in B(0,r)$ such that
\[
\mathbf{v}(x) = 0.
\]
\end{lemma}

\begin{proof}
\sketch If $\mathbf v$ had no zero in $B(0,r)$, the map $\mathbf w(x)=-r\mathbf v(x)/|\mathbf v(x)|$ would map the ball continuously to its boundary. Brouwer gives a fixed point $z$, but then $|z|=r$ and $r^2=\mathbf w(z)\cdot z\leq0$, a contradiction.
\end{proof}

\begin{theorem}[Construction of approximate solutions]
\label{thm:monotonicity-approximate-solutions}
\uses{lem:zeros-of-a-vector-field}
\dcref{ch9:9.1-thm-1}
\group{Nonlinear Elliptic PDE}

For each integer $m = 1, \ldots$, there exists a function $u_m$ of the form (6) satisfying the identities (7).
\end{theorem}

\begin{proof}
\sketch Define a continuous map on $\R^m$ whose components are the residuals in (7). Coercivity (5), Poincare's inequality, and the linear Dirichlet estimate imply that its scalar product with $d$ is nonnegative on a sufficiently large sphere. The zero lemma supplies coefficients $d_m^k$ with vanishing residual.
\end{proof}

\begin{theorem}[Energy estimates]
\label{thm:monotonicity-energy-estimates}
\uses{thm:monotonicity-approximate-solutions, thm:poincares-inequality-connected}
\dcref{ch9:9.1-thm-2}
\group{Energy Estimates}

There exists a constant $C$, depending only on $U$ and $\mathbf{a}$, such that
\[
\|u_m\|_{H^1_0(U)} \leq C(1 + \|f\|_{L^2(U)})
\tag{11}
\]
for $m = 1, 2, \ldots$.
\end{theorem}

\begin{proof}
\sketch Multiply (7) by $d_m^k$ and sum over $k$. Coercivity bounds $\|Du_m\|_2^2$ by a constant plus $\int_Ufu_m$; Cauchy's inequality and Poincare's inequality absorb the $u_m$ term and give (11).
\end{proof}

\begin{theorem}[Existence of weak solution]
\label{thm:monotonicity-existence-weak-solution}
\uses{thm:monotonicity-approximate-solutions, thm:monotonicity-energy-estimates, def:monotone-vector-field}
\dcref{ch9:9.1-thm-3}
\group{Nonlinear Elliptic PDE}

There exists a weak solution of the nonlinear boundary-value problem (1).
\end{theorem}

\begin{proof}
\sketch The energy estimate gives $u_m\rightharpoonup u$ in $H_0^1$ and, by (4), $\mathbf a(Du_m)\rightharpoonup\xi$ in $L^2$. Passing to the Galerkin identities gives the weak equation with $\xi$. Minty's monotonicity inequality, tested with $w=u-\lambda v$ and followed by $\lambda\downarrow0$, identifies $\xi=\mathbf a(Du)$ weakly and yields (12).
\end{proof}

\begin{remark}[Method of Browder and Minty]
\label{rem:browder-minty-method}
\uses{thm:monotonicity-existence-weak-solution}
\dcref{ch9:9.1-rem-1}
\group{Nonlinear Elliptic PDE}

The \textit{method of Browder and Minty} uses monotonicity to identify weak limits of nonlinear terms.
\end{remark}

Assume additionally that, for some $\theta>0$,
\[
(\mathbf a(p)-\mathbf a(q))\cdot(p-q)\geq\theta|p-q|^2
\quad(p,q\in\R^n).
\tag{20}
\]

\begin{theorem}[Uniqueness of weak solution]
\label{thm:monotonicity-uniqueness-weak-solution}
\uses{thm:monotonicity-existence-weak-solution}
\dcref{ch9:9.1-thm-4}
\group{Nonlinear Elliptic PDE}

Assume the strict monotonicity property (20) holds. Then there exists only one weak solution of (1).
\end{theorem}

\begin{proof}
\sketch Subtract the weak identities for two solutions and test with their difference. Strict monotonicity forces the $L^2$ norm of the gradient difference to vanish; the common zero trace then gives equality.
\end{proof}

\begin{remark}[$H^2$ regularity under strict monotonicity]
\label{rem:monotonicity-h2-regularity}
\uses{thm:monotonicity-uniqueness-weak-solution, thm:interior-h2-regularity}
\dcref{ch9:9.1-rem-2}
\group{Elliptic Regularity}

Under (20), the weak solution belongs to $H^2(U)$ and satisfies
\[
-\Div\,\mathbf{a}(Du) = f \quad \text{a.e. in } U.
\]
Indeed, applying (20) to $p=q+h\xi$ and sending $h\to0$ gives
\[
\sum_{i,j=1}^{n} a^i_{p_j}(q)\xi_i\xi_j \geq \theta|\xi|^2 \quad (q, \xi \in \R^n).
\tag{21}
\]
Thus the linearized coefficient matrix is uniformly elliptic, and the difference-quotient regularity argument applies (cf.\ \cref{thm:interior-h2-regularity}).
\end{remark}

\section{Fixed Point Methods}
\label{sec:fixed-point-methods}

\subsection{Banach's Fixed Point Theorem}
\label{sec:banachs-fixed-point-theorem}

Let $X$ be a Banach space.

\begin{theorem}[Banach's Fixed Point Theorem]
\label{thm:banach-fixed-point-theorem}
\dcref{ch9:9.2-thm-1}
\group{Fixed Point Theorems}
Assume
\[
A : X \to X
\]
is a nonlinear mapping, and suppose that
\[
\|A[u] - A[\bar{u}]\| \leq \gamma \|u - \bar{u}\| \quad (u, \bar{u} \in X)
\tag{1}
\]
for some constant $\gamma < 1$. Then $A$ has a unique fixed point.
\end{theorem}

\begin{proof}
\sketch Starting from $u_0\in X$, define $u_{k+1}=A[u_k]$. The contraction estimate bounds successive differences by a geometric series, so $(u_k)$ is Cauchy and converges to a fixed point. Applying the same estimate to two fixed points gives uniqueness.
\end{proof}

\begin{definition}[Strict contraction]
\label{def:strict-contraction}
\dcref{ch9:9.2.1-def-1}
\group{Fixed Point Theorems}
We say that $A$ is a strict contraction if (1) holds.
\end{definition}

\begin{example}[Reaction-diffusion equations]
\label{ex:reaction-diffusion-banach-fixed-point}
\uses{def:sobolev-space-time-banach}
\dcref{ch9:9.2-ex-1}
\group{Fixed Point Theorems}

Consider the \textit{reaction-diffusion system}
\[
\begin{cases}
\mathbf{u}_t - \Delta \mathbf{u} = f(\mathbf{u}) & \text{in } U_T \\
\mathbf{u} = 0 & \text{on } \partial U \times [0,T] \\
\mathbf{u} = \mathbf{g} & \text{on } U \times \{t = 0\}.
\end{cases}
\tag{2}
\]
Here $U_T=U\times(0,T]$, where $U\subset\R^n$ is bounded with smooth boundary and $T>0$; $\mathbf u=(u^1,\ldots,u^m)$ and $\mathbf g=(g^1,\ldots,g^m)\in H_0^1(U;\R^m)$. Assume
\[
f : \R^m \to \R^m \text{ is Lipschitz continuous.}
\tag{3}
\]
\end{example}

The Lipschitz hypothesis implies
\[
|\mathbf f(\mathbf z)|\leq C(1+|\mathbf z|).
\tag{4}
\]
Write $B[\mathbf u,\mathbf v]=\int_U D\mathbf u:D\mathbf v\,dx$. The weak-solution class is
\[
\mathbf u\in L^2(0,T;H_0^1(U;\R^m)),
\qquad
\mathbf u'\in L^2(0,T;H^{-1}(U;\R^m)),
\tag{5}
\]
with
\[
\langle\mathbf u',\mathbf v\rangle+B[\mathbf u,\mathbf v]
=(\mathbf f(\mathbf u),\mathbf v)
\quad\text{a.e. }0\leq t\leq T
\tag{6}
\]
for $\mathbf v\in H_0^1(U;\R^m)$, and
\[
\mathbf u(0)=\mathbf g.
\tag{7}
\]

\begin{definition}[Weak solution of the reaction-diffusion system]
\label{def:weak-solution-reaction-diffusion-system}
\uses{ex:reaction-diffusion-banach-fixed-point, thm:mappings-into-better-spaces}
\dcref{ch9:9.2.1-def-2}
\group{Fixed Point Theorems}

A function $\mathbf{u}$ satisfying (5) is a \textit{weak solution} of (2) provided (6) and (7) hold.
\end{definition}

\begin{theorem}[Existence for the reaction-diffusion system]
\label{thm:reaction-diffusion-banach-existence}
\uses{thm:banach-fixed-point-theorem, def:strict-contraction, def:weak-solution-reaction-diffusion-system}
\dcref{ch9:9.2-thm-2}
\group{Fixed Point Theorems}

There exists a unique weak solution of (2).
\end{theorem}

\begin{proof}
\sketch On $X=C([0,T];L^2(U;\R^m))$, map $\mathbf u$ to the solution of the linear heat system forced by $\mathbf f(\mathbf u)$. The parabolic energy estimate and Lipschitz continuity give $\|A\mathbf u-A\bar{\mathbf u}\|\leq(CT)^{1/2}\|\mathbf u-\bar{\mathbf u}\|$. Banach yields a solution on short intervals, iteration covers $[0,T]$, and the same energy inequality with Gronwall gives uniqueness.
\end{proof}

\begin{remark}[Chemical interpretation of the reaction-diffusion system]
\label{rem:reaction-diffusion-chemical-interpretation}
\uses{thm:reaction-diffusion-banach-existence}
\dcref{ch9:9.2.1-rem-1}
\group{Fixed Point Theorems}

In reaction-diffusion models, $\Delta\mathbf u$ represents diffusion and $\mathbf f(\mathbf u)$ represents reactions. Polynomial reaction laws need not be globally Lipschitz and may produce finite-time blow-up.
\end{remark}

\subsection{Schauder's, Schaefer's Fixed Point Theorems}
\label{sec:schauder-schaefer-fixed-point-theorems}

\begin{theorem}[Schauder's Fixed Point Theorem]
\label{thm:schauders-fixed-point-theorem}
\uses{thm:brouwer-fixed-point-theorem}
\dcref{ch9:9.2-thm-3}
\group{Fixed Point Theorems}

Suppose $K \subset X$ is compact and convex, and assume also
\[
A : K \to K
\]
is continuous. Then $A$ has a fixed point in $K$.
\end{theorem}

\begin{proof}
\sketch Cover the compact convex set by finitely many $\epsilon$-balls and map it continuously to the finite-dimensional convex hull of their centers with an error at most $\epsilon$. Brouwer gives a fixed point of the finite-dimensional approximation; compactness and continuity identify a limiting fixed point of $A$.
\end{proof}

\begin{definition}[Compact mapping]
\label{def:compact-mapping-banach-space}
\dcref{ch9:9.2.2-def-1}
\group{Fixed Point Theorems}
A nonlinear mapping $A : X \to X$ is called \textit{compact} provided for each bounded sequence $\{u_k\}_{k=1}^{\infty}$ the sequence $\{A[u_k]\}_{k=1}^{\infty}$ is precompact; that is, there exists a subsequence $\{u_{k_j}\}_{j=1}^{\infty}$ such that $\{A[u_{k_j}]\}_{j=1}^{\infty}$ converges in $X$.
\end{definition}

\begin{theorem}[Schaefer's Fixed Point Theorem]
\label{thm:schaefers-fixed-point-theorem}
\uses{thm:schauders-fixed-point-theorem, def:compact-mapping-banach-space}
\dcref{ch9:9.2-thm-4}
\group{Fixed Point Theorems}

Suppose
\[
A : X \to X
\]
is a continuous and compact mapping. Assume further that the set
\[
\{u \in X \mid u = \lambda A[u] \text{ for some } 0 \le \lambda \le 1\}
\]
is bounded. Then $A$ has a fixed point.
\end{theorem}

\begin{proof}
\sketch Choose $M$ larger than every solution of $u=\lambda A[u]$ and radially truncate $A$ to a compact map of $B(0,M)$ into itself. Schauder gives a fixed point of the truncation. If truncation were active there, the point would violate the defining a priori bound, so it is a fixed point of $A$.
\end{proof}

\begin{remark}[The method of a priori estimates]
\label{rem:a-priori-estimates-method}
\uses{thm:schaefers-fixed-point-theorem}
\dcref{ch9:9.2.2-rem-1}
\group{Fixed Point Theorems}

Bounding every solution of $u=\lambda A[u]$ uniformly for $0\leq\lambda\leq1$ yields a fixed point of $A$. This is the method of \textit{a priori estimates}.
\end{remark}

\begin{example}[A quasilinear elliptic PDE]
\label{ex:quasilinear-elliptic-schaefer}
\dcref{ch9:9.2-ex-2}
\group{Fixed Point Theorems}
Consider the semilinear boundary-value problem
\[
\begin{cases}
-\Delta u + b(Du) + \mu u = 0 & \text{in } U \\
u = 0 & \text{on } \partial U,
\end{cases}
\tag{21}
\]
where $U$ is bounded with smooth boundary and $b:\R^n\to\R$ is smooth and Lipschitz, hence
\[
|b(p)| \le C(|p| + 1)
\tag{22}
\]
for some constant $C$ and all $p \in \R^n$.
\end{example}

\begin{theorem}[Existence for the quasilinear elliptic PDE]
\label{thm:schaefer-quasilinear-existence}
\uses{thm:schaefers-fixed-point-theorem, def:compact-mapping-banach-space, thm:second-derivatives-minimizers}
\dcref{ch9:9.2-thm-5}
\group{Fixed Point Theorems}

If $\mu > 0$ is sufficiently large, there exists a function $u \in H^2(U) \cap H^1_0(U)$ solving the boundary-value problem (21).
\end{theorem}

\begin{proof}
\sketch Map $u\in H_0^1$ to the solution $w$ of $-\Delta w+\mu w=-b(Du)$. Elliptic $H^2$ estimates and compact embedding make this map continuous and compact. Testing the equation for $u=\lambda A[u]$ by $u$ gives a uniform $H_0^1$ bound when $\mu$ is large, so Schaefer supplies a fixed point in $H^2\cap H_0^1$.
\end{proof}

\begin{remark}[Remark and warning]
\label{rem:iteration-scheme-warning}
\uses{thm:schaefer-quasilinear-existence}
\dcref{ch9:9.2.2-rem-2}
\group{Fixed Point Theorems}

The iteration
\[
\begin{cases}
-\Delta u^{k+1} + \mu u^{k+1} = -b(Du^k) & \text{in } U \\
u^{k+1} = 0 & \text{on } \partial U
\end{cases} \quad (k = 0, 1, \ldots).
\]
is not guaranteed to converge to a solution of (21): Schauder's and Schaefer's theorems assert existence of fixed points, not convergence of iterates.
\end{remark}

\section{Method of Subsolutions and Supersolutions}
\label{sec:subsolutions-supersolutions}

Let $f:\R\to\R$ be smooth with $|f'(z)|\leq C$, and consider
\[
\begin{cases}
-\Delta u=f(u) & \text{in }U,\\
u=0 & \text{on }\partial U.
\end{cases}
\tag{1}
\]

\begin{definition}[Weak sub- and supersolutions]
\label{def:weak-sub-supersolution}
\dcref{ch9:9.3-def-1}
\group{Sub- and Supersolutions}
\begin{enumerate}
\item[(i)] We say that $\bar{u} \in H^1(U)$ is a \textit{weak supersolution} of problem (1) if
\[
\int_U D\bar{u} \cdot Dv \, dx \geq \int_U f(\bar{u})v \, dx
\tag{3}
\]
for each $v \in H^1_0(U), v \geq 0$ a.e.
\item[(ii)] Similarly, $\underline{u} \in H^1(U)$ is a \textit{weak subsolution} provided
\[
\int_U D\underline{u} \cdot Dv \, dx \leq \int_U f(\underline{u})v \, dx
\tag{4}
\]
for each $v \in H^1_0(U), v \geq 0$ a.e.
\item[(iii)] We say $u \in H^1_0(U)$ is a \textit{weak solution} of (1) if
\[
\int_U Du \cdot Dv \, dx = \int_U f(u)v \, dx
\]
for each $v \in H^1_0(U)$.
\end{enumerate}
\end{definition}

\begin{remark}[Classical interpretation]
\label{rem:sub-supersolution-c2-case}
\uses{def:weak-sub-supersolution}
\dcref{ch9:9.3-rem-1}
\group{Sub- and Supersolutions}

If $\bar{u}, \underline{u} \in C^2(U)$, then from (3) and (4) it follows that
\[
-\Delta \bar{u} \geq f(\bar{u}), \quad -\Delta \underline{u} \leq f(\underline{u}) \quad \text{in } U.
\]
\end{remark}

\begin{theorem}[Existence of a solution between sub- and supersolutions]
\label{thm:existence-between-sub-supersolutions}
\uses{def:weak-sub-supersolution, thm:first-existence-weak-solutions}
\dcref{ch9:9.3-thm-1}
\group{Sub- and Supersolutions}

Assume there exist a weak supersolution $\bar{u}$ and a weak subsolution $\underline{u}$ of (1), satisfying
\[
\underline{u} \leq 0, \quad \bar{u} \geq 0 \quad \text{on } \partial U \text{ in the trace sense}, \quad \underline{u} \leq \bar{u} \text{ a.e. in } U.
\tag{5}
\]
Then there exists a weak solution $u$ of (1), such that
\[
\underline{u} \leq u \leq \bar{u} \quad \text{a.e. in } U.
\]
\end{theorem}

\begin{proof}
\sketch Choose $\lambda$ so that $z\mapsto f(z)+\lambda z$ is nondecreasing and iterate the linear problem $-\Delta u_{k+1}+\lambda u_{k+1}=f(u_k)+\lambda u_k$ from $u_0=\underline u$. Testing differences with positive parts proves $\underline u\leq u_k\leq u_{k+1}\leq\bar u$. Dominated and weak compactness give a limit, and passage to the weak identity makes it a solution between the barriers.
\end{proof}

\section{Nonexistence}
\label{sec:nonexistence}

\subsection{Blow-up}
\label{sec:blow-up}

Consider
\[
\begin{cases}
u_t-\Delta u=u^2 & \text{in }U_T,\\
u=0 & \text{on }\partial U\times(0,T),\\
u=g & \text{on }U\times\{t=0\}.
\end{cases}
\tag{1}
\]
Let $\lambda_1>0$ be the principal eigenvalue of $-\Delta$ on $H_0^1(U)$ and choose $w_1$ with
\[
\begin{cases}
-\Delta w_1=\lambda_1w_1 & \text{in }U,\\
w_1=0 & \text{on }\partial U,
\end{cases}
\qquad
w_1>0,
\qquad
\int_Uw_1\,dx=1.
\tag{3}
\]
For a smooth positive solution define
\[
\eta(t):=\int_Uu(x,t)w_1(x)\,dx.
\]
Then integration by parts and Cauchy's inequality give
\[
\dot\eta(t)\geq-\lambda_1\eta(t)+\eta(t)^2.
\tag{6}
\]

\begin{theorem}[Blow-up for large initial data]
\label{thm:blowup-large-data}
\uses{def:principal-eigenvalue}
\dcref{ch9:9.4.1-thm-1}
\group{Nonexistence}

Let $u$ be a smooth solution of (1) with $g\geq0$, $g\not\equiv0$, let $w_1$ satisfy (3), and define $\eta(t)=\int_Uu(x,t)w_1(x)\,dx$. If
\[
\eta(0) = \int_U gw_1 \, dx > \lambda_1,
\]
set
\[
t^*:=-\frac1{\lambda_1}\log\left(\frac{\eta(0)-\lambda_1}{\eta(0)}\right).
\tag{10}
\]
Then either smoothness fails before $t^*$ or
\[
\lim_{t \to t_*} \int_U u(x,t)w_1(x) \, dx = \infty
\]
for some $0<t_*\leq t^*$. In particular, no smooth solution exists on every interval $[0,T]$ with $T\geq t^*$.
\end{theorem}

\begin{proof}
\sketch Set $\xi(t)=e^{\lambda_1t}\eta(t)$. Inequality (6) gives $\dot\xi\geq e^{-\lambda_1t}\xi^2$, whose integrated reciprocal inequality has a denominator vanishing at $t^*$. Thus $\eta$ becomes unbounded no later than $t^*$ unless the assumed smooth solution ceases to exist first.
\end{proof}

\subsection{Derrick--Pohozaev Identity}
\label{sec:derrick-pohozaev-identity}

For $n>2$, consider
\[
\begin{cases}
-\Delta u=|u|^{p-1}u & \text{in }U,\\
u=0 & \text{on }\partial U,
\end{cases}
\tag{11}
\]
in the supercritical range
\[
\frac{n+2}{n-2}<p<\infty.
\tag{13}
\]

\begin{definition}[Star-shaped domain]
\label{def:star-shaped-domain}
\dcref{ch9:9.4.2-def-1}
\group{Nonexistence}
An open set $U$ is called \textit{star-shaped with respect to} $0$ provided for each $x \in \overline{U}$, the line segment
\[
\{\lambda x \mid 0 \leq \lambda \leq 1\}
\]
lies in $\overline{U}$.
\end{definition}

\begin{lemma}[Normals to a star-shaped region]
\label{lem:normals-star-shaped-region}
\uses{def:star-shaped-domain}
\dcref{ch9:9.4.2-lem-1}
\group{Nonexistence}

Assume $\partial U$ is $C^1$ and $U$ is star-shaped with respect to 0. Then
\[
x \cdot \nu(x) \geq 0 \quad \text{for all } x \in \partial U,
\]
where $\nu$ denotes the unit outward normal.
\end{lemma}

\begin{proof}
\sketch For $x\in\partial U$, star-shapedness puts $\lambda x$ in $\overline U$ for $0<\lambda<1$. The $C^1$ boundary tangent-cone inequality applied as $\lambda\uparrow1$ gives $\nu(x)\cdot x/|x|\geq0$.
\end{proof}

\begin{theorem}[Nonexistence of nontrivial solution]
\label{thm:derrick-pohozaev-nonexistence}
\uses{def:star-shaped-domain, lem:normals-star-shaped-region}
\dcref{ch9:9.4.2-thm-1}
\group{Nonexistence}

Assume $u \in C^2(\overline{U})$ is a solution of problem (11), and the exponent $p$ satisfies inequality (13). Suppose further $U$ is star-shaped with respect to 0, and $\partial U$ is $C^1$. Then
\[
u \equiv 0 \quad \text{within } U.
\]
\end{theorem}

\begin{proof}
\sketch Multiply (11) by $x\cdot Du$ and integrate by parts to obtain the Pohozaev identity. Star-shapedness makes its boundary term nonnegative. Combining the resulting inequality with the identity obtained by testing (11) with $u$ forces $p\leq(n+2)/(n-2)$ unless $u=0$, contradicting (13).
\end{proof}

\begin{remark}[Derrick--Pohozaev identity]
\label{rem:derrick-pohozaev-identity-name}
\uses{thm:derrick-pohozaev-nonexistence}
\dcref{ch9:9.4.2-rem-1}
\group{Nonexistence}

The \textit{Derrick--Pohozaev identity} is
\[
\frac{n-2}{2}\int_U|Du|^2\,dx
+\frac12\int_{\partial U}|Du|^2(\nu\cdot x)\,dS
=\frac{n}{p+1}\int_U|u|^{p+1}\,dx.
\tag{18}
\]
\end{remark}

\section{Geometric Properties of Solutions}
\label{sec:geometric-properties-of-solutions}

\subsection{Star-shaped Level Sets}
\label{sec:star-shaped-level-sets}

Let $V\Subset W$ be open sets star-shaped with respect to $0$, set
$U=W\setminus\overline V$, $\Gamma_0=\partial W$, and $\Gamma_1=\partial V$, and let $u$ solve
\[
\begin{cases}
-\Delta u=0 & \text{in }U,\\
u=1 & \text{on }\Gamma_1,\\
u=0 & \text{on }\Gamma_0.
\end{cases}
\tag{1}
\]
Then $0<u<1$ in $U$.

\begin{theorem}[Star-shaped level sets]
\label{thm:star-shaped-level-sets}
\uses{def:star-shaped-domain, lem:normals-star-shaped-region, thm:strong-maximum-principle-laplace}
\dcref{ch9:9.5.1-thm-1}
\group{Symmetry of Solutions}

For each $0 < \lambda < 1$ the level set
\[
\Gamma_\lambda := \{x \in U \mid u(x) = \lambda\}
\]
is a smooth surface and is the boundary of a set star-shaped with respect to $0$.
\end{theorem}

\begin{proof}
\sketch The scaling derivative $v(x)=Du(x)\cdot x$ is harmonic. Star-shapedness and the boundary data give $v\leq0$ on both boundary components, so the strong maximum principle yields $v<0$ in $U$. Hence $Du\neq0$, each level is smooth, and its outward normal satisfies $x\cdot\nu>0$; the radial segment criterion then makes the enclosed superlevel set star-shaped.
\end{proof}

\subsection{Radial Symmetry}
\label{sec:radial-symmetry}

Let $U=B^0(0,1)\subset\R^n$, let $f:\R\to\R$ be Lipschitz, and suppose
\[
\begin{cases}
-\Delta u=f(u) & \text{in }U,\\
u=0 & \text{on }\partial U,
\end{cases}
\tag{2}
\qquad
u>0\quad\text{in }U.
\tag{3}
\]

\subsubsection{Maximum Principles}
\label{sec:maximum-principles-radial-symmetry}

\begin{lemma}[A refinement of Hopf's Lemma]
\label{lem:hopf-lemma-refinement}
\uses{lem:hopfs-lemma, thm:strong-maximum-principle-c-zero}
\dcref{ch9:9.5.2-lem-1}
\group{Maximum Principles}

Suppose $V \subset \R^n$ is open, $v \in C^2(\overline{V})$, and $c \in L^\infty(V)$. Assume
\[
\begin{cases} -\Delta v + cv \geq 0 & \text{in } V \\ v \geq 0 & \text{in } V. \end{cases}
\tag{4}
\]
Suppose also $v \not\equiv 0$.
\begin{enumerate}
\item[(i)] If $x^0 \in \partial V$, $v(x^0) = 0$, and $V$ satisfies the interior ball condition at $x^0$, then
\[
\frac{\partial v}{\partial \nu}(x^0) < 0.
\tag{5}
\]
\item[(ii)] Furthermore,
\[
v > 0 \text{ in } V.
\tag{6}
\]
\end{enumerate}
No sign condition on the zeroth-order coefficient $c$ is required.
\end{lemma}

\begin{proof}
\sketch Set $w=e^{-\lambda x_1}v$ with $\lambda^2\geq\|c\|_\infty$. Then $w$ is a nonnegative supersolution of an elliptic operator with no zeroth-order term. The strong maximum principle gives $w>0$, and Hopf's lemma transfers the strict outward derivative inequality from $w$ to $v$.
\end{proof}

\begin{lemma}[Boundary estimates]
\label{lem:boundary-estimates-radial-symmetry}
\uses{lem:hopf-lemma-refinement}
\dcref{ch9:9.5.2-lem-2}
\group{Maximum Principles}

Let $u \in C^2(\overline{U})$ satisfy (2), (3). Then for each point $x^0 \in \partial U \cap \{x_n > 0\}$, either
\[
u_{x_n}(x^0) < 0
\tag{7}
\]
or else
\[
u_{x_n}(x^0) = 0, \quad u_{x_nx_n}(x^0) > 0.
\tag{8}
\]
In either case, $u$ is strictly decreasing as a function of $x_n$ near $x^0$.
\end{lemma}

\begin{proof}
\sketch If $f(0)\geq0$, apply the refined Hopf lemma to $u$ after writing $-\Delta u+c(x)u\geq0$. If $f(0)<0$ and the first normal derivative vanishes, differentiate the boundary identity twice and use the PDE to obtain $D^2u(x^0)=-f(0)\nu\otimes\nu$, giving the positive second derivative in (8).
\end{proof}

\subsubsection{Moving Planes}
\label{sec:moving-planes}

\begin{definition}[Moving plane notation]
\label{def:moving-plane-notation}
\dcref{ch9:9.5.2-def-1}
\group{Maximum Principles}
\begin{enumerate}
\item[(i)] If $0 \le \lambda \le 1$, define the plane
\[
P_\lambda := \{x \in \R^n \mid x_n = \lambda\}.
\]
\item[(ii)] Write $x_\lambda := (x_1, \ldots, x_{n-1}, 2\lambda - x_n)$ to denote the \textit{reflection} of $x$ in $P_\lambda$.
\item[(iii)] $E_\lambda := \{x \in U \mid \lambda < x_n < 1\}$.
\end{enumerate}
\end{definition}

\begin{theorem}[Radial symmetry]
\label{thm:radial-symmetry-positive-solutions}
\uses{lem:hopf-lemma-refinement, lem:boundary-estimates-radial-symmetry, def:moving-plane-notation}
\dcref{ch9:9.5.2-thm-2}
\group{Symmetry of Solutions}

Let $u \in C^2(\overline{U})$ solve (2), (3). Then $u$ is radial; that is,
\[
u(x) = v(r) \quad (r = |x|)
\]
for some strictly decreasing function $v : [0, 1] \to [0, \infty)$.
\end{theorem}

\begin{proof}
\sketch Reflect the cap $E_\lambda$ across $P_\lambda$ and compare $u(x_\lambda)$ with $u(x)$. Boundary estimates start the strict inequality near $\lambda=1$; the refined maximum principle and Hopf inequality prevent its first failure as the plane moves to zero. Reflection symmetry across every hyperplane through the origin yields radiality, and the strict comparison yields radial decrease.
\end{proof}

\section{Gradient Flows}
\label{sec:gradient-flows}

\subsection{Convex Functions on Hilbert Spaces}
\label{sec:convex-functions-hilbert-spaces}

Let $H$ be a real Hilbert space with inner product $(\cdot,\cdot)$ and norm $\|\cdot\|$.

\begin{definition}[Convex function on a Hilbert space]
\label{def:convex-function-hilbert-space}
\dcref{ch9:9.6.1-def-1}
\group{Convex Analysis}
A function
\[
I : H \to (-\infty, \infty]
\]
is \textit{convex} provided
\[
I[\tau u + (1-\tau)v] \le \tau I[u] + (1-\tau)I[v]
\]
for all $u, v \in H$ and each $0 \le \tau \le 1$.
\end{definition}

A function $I:H\to(-\infty,+\infty]$ is \textit{proper} if it is not identically $+\infty$, and its domain is
\[
D(I):=\{u\in H\mid I[u]<+\infty\}.
\]

\begin{definition}[Lower semicontinuous]
\label{def:lower-semicontinuous-hilbert-functional}
\dcref{ch9:9.6.1-def-2}
\group{Convex Analysis}
We say $I : H \to (-\infty, +\infty]$ is \textit{lower semicontinuous} if
\[
\begin{cases} u_k \to u \text{ in } H \text{ implies} \\ I[u] \le \liminf_{k \to \infty} I[u_k]. \end{cases}
\]
\end{definition}

\begin{definition}[Subdifferential]
\label{def:subdifferential-hilbert}
\uses{def:convex-function-hilbert-space}
\dcref{ch9:9.6.1-def-3}
\group{Convex Analysis}

Let $I : H \to (-\infty, +\infty]$ be convex and proper.
\begin{enumerate}
\item[(i)] For each $u \in H$, we write
\[
\partial I[u] := \{v \in H \mid I[w] \ge I[u] + (v, w-u) \text{ for all } w \in H\}.
\tag{1}
\]
The mapping $\partial I : H \to 2^H$ is the \textit{subdifferential} of $I$.
\item[(ii)] We say $u \in D(\partial I)$, the domain of $\partial I$, provided $\partial I[u] \ne \emptyset$.
\end{enumerate}
\end{definition}

\begin{theorem}[Properties of subdifferentials]
\label{thm:properties-of-subdifferentials}
\uses{def:subdifferential-hilbert, def:lower-semicontinuous-hilbert-functional, thm:weak-lower-semicontinuity-convex-lagrangian}
\dcref{ch9:9.6-thm-1}
\group{Convex Analysis}

Let $I : H \to (-\infty, +\infty]$ be convex, proper and lower semicontinuous. Then
\begin{enumerate}
\item[(i)] $D(\partial I) \subseteq D(I)$.
\item[(ii)] If $v \in \partial I[u]$ and $\tilde{v} \in \partial I[\tilde{u}]$, then
\[
(v - \tilde{v}, u - \tilde{u}) \geq 0 \quad \text{(monotonicity)}.
\]
\item[(iii)] $I[u] = \min_{w \in H} I[w]$ if and only if $0 \in \partial I[u]$.
\item[(iv)] For each $w \in H$ and $\lambda > 0$, the problem
\[
u + \lambda \partial I[u] \ni w
\]
has a unique solution $u \in D(\partial I)$.
\end{enumerate}
Equivalently, in (iv) there exist $u\in D(\partial I)$ and $v\in\partial I[u]$ such that
\[
u + \lambda v = w.
\]
\end{theorem}

\begin{proof}
\sketch The supporting inequalities immediately give (i), monotonicity, and the minimizer criterion. For (iv), minimize $J[u]=\frac12\|u\|^2+\lambda I[u]-(u,w)$. Convex lower semicontinuity implies weak lower semicontinuity, while a global affine lower bound for $I$ makes $J$ coercive. Its minimizer solves $u+\lambda\partial I[u]\ni w$, and monotonicity gives uniqueness.
\end{proof}

\begin{definition}[Nonlinear resolvent and Yosida approximation]
\label{def:nonlinear-resolvent-yosida}
\uses{thm:properties-of-subdifferentials}
\dcref{ch9:9.6.1-def-4}
\group{Convex Analysis}

\begin{enumerate}
\item[(1)] For each $\lambda > 0$ define the \textit{nonlinear resolvent} $J_\lambda : H \to D(\partial I)$ by setting
\[
J_\lambda[w] := u,
\]
where $u$ is the unique solution of
\[
u + \lambda \partial I[u] \ni w.
\]
\item[(2)] For each $\lambda > 0$ define the \textit{Yosida approximation} $A_\lambda : H \to H$ by
\[
A_\lambda[w] := \frac{w - J_\lambda[w]}{\lambda} \quad (w \in H).
\tag{7}
\]
\end{enumerate}
The operator $A_\lambda$ is the Yosida regularization of $A=\partial I$.
\end{definition}

\begin{theorem}[Properties of $J_\lambda$, $A_\lambda$]
\label{thm:properties-resolvent-yosida}
\uses{def:nonlinear-resolvent-yosida, thm:properties-of-subdifferentials}
\dcref{ch9:9.6-thm-2}
\group{Convex Analysis}

For each $\lambda > 0$ and $w, \tilde{w} \in H$, the following statements hold:
\begin{enumerate}
\item[(i)] $\|J_\lambda[w] - J_\lambda[\tilde{w}]\| \le \|w - \tilde{w}\|$,
\item[(ii)] $\|A_\lambda[w] - A_\lambda[\tilde{w}]\| \le \frac{2}{\lambda}\|w - \tilde{w}\|$,
\item[(iii)] $0 \le (w - \tilde{w}, A_\lambda[w] - A_\lambda[\tilde{w}])$,
\item[(iv)] $A_\lambda[w] \in \partial I[J_\lambda[w]]$.
\item[(v)] If $w \in D(\partial I)$, then
\[
\sup_{\lambda > 0} \|A_\lambda[w]\| \le |A^0[w]|,
\]
where $|A^0[w]| := \min_{z \in \partial I[w]} \|z\|$.
\item[(vi)] For each $w \in \overline{D(\partial I)}$,
\[
\lim_{\lambda \to 0} J_\lambda[w] = w.
\]
\end{enumerate}
\end{theorem}

\begin{proof}
\sketch Write $J_\lambda w+\lambda v=w$ and $J_\lambda\tilde w+\lambda\tilde v=\tilde w$ and use monotonicity of $\partial I$. This gives the resolvent contraction, the Lipschitz and monotonicity estimates for $A_\lambda$, and $A_\lambda w\in\partial I[J_\lambda w]$. Comparing with the minimum-norm element of $\partial I[w]$ yields (v); the identity $J_\lambda w-w=-\lambda A_\lambda w$, followed by density, yields (vi).
\end{proof}

\subsection{Subdifferentials and Nonlinear Semigroups}
\label{sec:subdifferentials-nonlinear-semigroups}

Let $I:H\to(-\infty,+\infty]$ be convex, proper, and lower semicontinuous, and assume
\[
\overline{D(\partial I)}=H.
\tag{8}
\]
For the gradient flow
\[
\begin{cases}
\mathbf u'(t)+\partial I[\mathbf u(t)]\ni0 & (t\geq0),\\
\mathbf u(0)=u,
\end{cases}
\tag{9}
\]
write
\[
\mathbf u(t)=S(t)u.
\tag{10}
\]
The semigroup identities and continuity condition are
\[
S(0)u=u,
\tag{11}
\]
\[
S(t+s)u=S(t)S(s)u,
\tag{12}
\]
and
\[
t\longmapsto S(t)u\text{ is continuous from }[0,\infty)\text{ to }H.
\tag{13}
\]

\begin{remark}[Nonlinear semigroup notation]
\label{rem:nonlinear-semigroup-notation}
\dcref{ch9:9.6.2-rem-1}
\group{Nonlinear Semigroups}
For each $t\geq0$, the map $u\mapsto S(t)u$ in (10) is generally nonlinear.
\end{remark}

\begin{definition}[Nonlinear (contraction) semigroup]
\label{def:nonlinear-contraction-semigroup}
\dcref{ch9:9.6.2-def-1}
\group{Nonlinear Semigroups}
\begin{enumerate}
\item[(i)] A family $\{S(t)\}_{t \geq 0}$ of nonlinear operators mapping $H$ into $H$ is called a \textit{nonlinear semigroup} if conditions (11)--(13) are satisfied.
\item[(ii)] We say $\{S(t)\}_{t \geq 0}$ is a \textit{contraction semigroup} if in addition
\[
\|S(t)u - S(t)\hat{u}\| \leq \|u - \hat{u}\| \quad (t \geq 0, u, \hat{u} \in H).
\tag{14}
\]
\end{enumerate}
\end{definition}

\begin{theorem}[Solution of gradient flow]
\label{thm:solution-of-gradient-flow}
\uses{thm:properties-resolvent-yosida, def:nonlinear-contraction-semigroup}
\dcref{ch9:9.6-thm-3}
\group{Nonlinear Semigroups}

For each $u \in D(\partial I)$ there exists a unique function
\[
\mathbf{u} \in C([0, \infty); H), \text{ with } \mathbf{u}' \in L^\infty(0, \infty; H),
\tag{16}
\]
such that
\begin{enumerate}
\item[(i)] $\mathbf{u}(0) = u$,
\item[(ii)] $\mathbf{u}(t) \in D(\partial I)$ for each $t > 0$,
\end{enumerate}
and
\begin{enumerate}
\item[(iii)] $\mathbf{u}'(t) \in -\partial I[\mathbf{u}(t)]$ for a.e. $t \geq 0$.
\end{enumerate}
\end{theorem}

\begin{proof}
\sketch Solve the regularized ODE $\mathbf u_\lambda'+A_\lambda(\mathbf u_\lambda)=0$. Monotonicity gives contraction, a uniform derivative bound, and a Cauchy estimate as $\lambda\downarrow0$, producing a continuous limit with bounded derivative. Passing the subdifferential inequality to the limit shows $-\mathbf u'(t)\in\partial I[\mathbf u(t)]$ and closedness gives domain membership at every time. Monotonicity applied to two solutions yields uniqueness.
\end{proof}

\begin{remark}[Extension to $\overline{D(\partial I)}$]
\label{rem:gradient-flow-semigroup-extension}
\uses{thm:solution-of-gradient-flow}
\dcref{ch9:9.6.2-rem-2}
\group{Nonlinear Semigroups}

For $\mathbf u,\mathbf v\in D(\partial I)$, the limiting flow satisfies
\[
\|S(t)\mathbf{u} - S(t)\mathbf{v}\| \leq \|\mathbf{u} - \mathbf{v}\| \quad (t \geq 0)
\]
and therefore extends uniquely to a nonlinear contraction semigroup on $\overline{D(\partial I)}$. Under (8), this space is $H$.
\end{remark}

\subsection{Applications}
\label{sec:gradient-flow-applications}

Let $U\subset\R^n$ be bounded with smooth boundary, take $H=L^2(U)$, and define
\[
I[u]:=
\begin{cases}
\displaystyle\int_U L(Du)\,dx & \text{if }u\in H_0^1(U),\\
+\infty & \text{otherwise},
\end{cases}
\tag{27}
\]
where $L:\R^n\to\R$ is smooth and convex, with
\[
|D^2L(p)|\leq C
\tag{28}
\]
and
\[
\sum_{i,j=1}^nL_{p_ip_j}(p)\xi_i\xi_j\geq\theta|\xi|^2
\quad(p,\xi\in\R^n)
\tag{29}
\]
for constants $C,\theta>0$.

\begin{theorem}[Characterization of $\partial I$]
\label{thm:characterization-subdifferential-uniformly-convex-lagrangian}
\uses{def:convex-function-hilbert-space, def:subdifferential-hilbert, thm:weak-lower-semicontinuity-convex-lagrangian, def:monotone-vector-field, thm:second-derivatives-minimizers}
\dcref{ch9:9.6-thm-4}
\group{Convex Analysis}

\begin{enumerate}
\item[(i)] $I : H \to (-\infty, +\infty]$ is convex, proper and lower semicontinuous.
\item[(ii)] $D(\partial I) = H^2(U) \cap H_0^1(U)$.
\item[(iii)] If $u \in D(\partial I)$, then $\partial I$ is single-valued and
\[
\partial I[u] = -\sum_{i=1}^n (L_{p_i}(Du))_{x_i} \quad \text{a.e.}
\]
\end{enumerate}
\end{theorem}

\begin{proof}
\sketch Convex weak lower semicontinuity gives (i). For $A[u]=-\Div D_pL(Du)$ on $H^2\cap H_0^1$, the supporting inequality for $L$ shows $A\subseteq\partial I$. Uniform convexity and elliptic regularity solve $u+A[u]=f$ for every $f\in L^2$, so $\operatorname{Id}+A$ is surjective; resolvent uniqueness then forces $A=\partial I$, proving (ii)--(iii).
\end{proof}

\begin{example}[Quasilinear parabolic PDE via gradient flows]
\label{ex:quasilinear-parabolic-gradient-flow}
\uses{thm:characterization-subdifferential-uniformly-convex-lagrangian, thm:solution-of-gradient-flow}
\dcref{ch9:9.6.3-ex-1}
\group{Nonlinear Semigroups}

Consider
\[
\begin{cases}
u_t - \sum_{i=1}^n (L_{p_i}(Du))_{x_i} = 0 & \text{in } U \times (0, \infty) \\
u = 0 & \text{on } \partial U \times (0, \infty) \\
u = g & \text{on } U \times \{t = 0\},
\end{cases}
\tag{33}
\]
where $g \in L^2(U)$. By the characterization of $\partial I$, this is the abstract flow
\[
\begin{cases} u'(t) = -\partial I[u(t)] & (t \geq 0) \\ u(0) = g. \end{cases}
\tag{34}
\]
If $g \in H^2(U) \cap H^1_0(U)$, the gradient-flow theorem gives a unique function
\[
u \in C([0,\infty); L^2(U)), \text{ with } u' \in L^\infty((0,\infty); L^2(U)),
\]
which is a weak solution of (33). The estimate
\[
\|u(t)\|_{H^2(U)} \leq C\|u'(t)\|_{L^2(U)},
\]
also gives $u \in L^\infty((0,\infty); H^2(U) \cap H^1_0(U))$.
\end{example}

\section{References}
\label{sec:references-ch9}
\textbf{Monotonicity and fixed points.} \cite{Lions1969}, \cite{Zeidler1985}, \cite{Evans1990}, \cite{Zeidler1985}, \cite{GilbargTrudinger1983}, \cite{GilbargTrudinger1983}.
\textbf{Subsolutions, nonexistence, and symmetry.} \cite{Smoller1983}, \cite{Payne1975}, \cite{GidasNiNirenberg1979}.
\textbf{Nonlinear semigroups.} \cite{Brezis1973}, \cite{Barbu1976}, \cite{Zeidler1985}.
